%%% ============================================================
%%% Logarithmic Quintessence: A Dimensionless Scalar Dark Energy
%%% Model with an Analytic Vacuum
%%%
%%% B.R. Sanders, M. Gish (2026)
%%% Target venue: JCAP (Journal of Cosmology and Astroparticle Physics)
%%% MSC: 83F05, 83C56, 85A40
%%% PACS: 95.36.+x (dark energy); 98.80.Es (observational cosmology);
%%%       98.80.Cq (cosmological perturbations)
%%%
%%% DOI for verification scripts: 10.5281/zenodo.18852047
%%% ============================================================

\documentclass[11pt,reqno]{amsart}

\usepackage{amsmath,amssymb,amsthm,mathtools}
\usepackage[margin=1in]{geometry}
\usepackage[colorlinks=true,linkcolor=blue,citecolor=blue,urlcolor=blue]{hyperref}
\usepackage{microtype}
\usepackage{array}

%% -------- theorem environments --------
\theoremstyle{plain}
\newtheorem{theorem}{Theorem}
\newtheorem{proposition}[theorem]{Proposition}
\newtheorem{corollary}[theorem]{Corollary}
\newtheorem{lemma}[theorem]{Lemma}

\theoremstyle{definition}
\newtheorem{definition}[theorem]{Definition}

\theoremstyle{remark}
\newtheorem*{remark}{Remark}

%% -------- macros --------
\newcommand{\Pl}{\mathrm{Pl}}
\newcommand{\SM}{\mathrm{SM}}
\newcommand{\Gibbs}{\mathrm{Gibbs}}
\DeclareMathOperator{\diag}{diag}

%% -------- title matter --------
\title[Logarithmic Quintessence]{Logarithmic Quintessence:\\
A Dimensionless Scalar Dark Energy Model with\\
an Analytic Vacuum}

\author{B.~R.~Sanders}
\address{7Site LLC, Hot Springs, Arkansas, USA}
\email{brayden@7site.co}

\author{M.~Gish}
\address{Independent Researcher, Hot Springs, Arkansas, USA}
\email{monica.gish1992@gmail.com}

\date{\today}

\keywords{quintessence, dark energy, scalar field, logarithmic potential,
freezing equation of state, entropy maximum, DESI, information-theoretic cosmology}

\subjclass[2020]{83F05, 83C56, 85A40}

\begin{document}

\begin{abstract}
We study a quintessence dark-energy model in which a real, positive,
dimensionless scalar field $\Xi(x)$ is minimally coupled to gravity with
self-interaction $V(\Xi) = \Lambda^{4}\,\Xi\log\Xi$, where $\Lambda$ is a
mass scale that sets the potential energy density. The potential admits an
analytic minimum at $\Xi_0 = e^{-1}$. As is generic for any dimensionless
field whose coupling enters as an overall prefactor, the location of this
minimum is independent of $\Lambda$; what is specific to the choice
$V \propto \Xi\log\Xi$ is the value $\Xi_0 = e^{-1}$ itself. Expanding about
$\Xi_0$ gives a stable massive scalar with
$m_\Xi^{2} = \Lambda^{4} e / M_\Pl^{2}$, which places $\Lambda$ numerically
near the observed dark-energy scale when $m_\Xi \sim H_0$, giving
$\Lambda \approx 1.7$~meV. In a spatially flat FRW background, the field
equation drives a freezing quintessence history with asymptotic limit
$w \to -1$ as the field settles at the vacuum. We provide an illustrative
background-level numerical integration showing that the model can reproduce
$(w_0, w_a)$ values close to the DESI~2024 DR1 published Gaussian summary
on the CPL parametrization for appropriately chosen initial conditions;
the $\chi^{2}$ values reported here quantify proximity to the published
$(w_0, w_a)$ marginal Gaussian and are not derived from the underlying
joint BAO + CMB + SN likelihood, which is deferred to a companion paper.
The minimal action couples to matter only gravitationally, so no
observable fifth force is expected from the minimal theory. The functional
form $-\Xi\log\Xi$ coincides with the per-bin Gibbs integrand on a
normalized probability distribution and with the
Bialynicki-Birula--Mycielski nonlinearity in nonlinear quantum mechanics
\cite{BialynickiBirulaMycielski1976}; both connections are structural
rhymes (the present $\Xi$ is unbounded and unnormalized, not a probability
density), cited as motivation for studying this functional form rather
than as derivations. Falsifiability is anchored in three Stage-IV
predictions: monotone freezing $w(z)$, asymptotic limit $w \to -1$, and a
specific two-parameter $w(z)$ profile.
\end{abstract}

\maketitle

\section{Introduction}

The nature of dark energy is the central open question in observational
cosmology. Recent results from the Dark Energy Spectroscopic Instrument
\cite{DESI2024BAO,DESI2024VI} suggest a preference for dynamical dark energy
over the cosmological constant at the $2.8$--$4.2\sigma$ level in joint
analyses combining DESI BAO with CMB and supernova data, in the
Chevallier-Polarski-Linder \cite{CPL2001} parametrization
$w(z) = w_0 + w_a z/(1+z)$. The DR1 BAO data alone are consistent with flat
$\Lambda$CDM; the dynamical-dark-energy preference appears in combined
analyses. These results have revived interest in scalar-field models of
quintessence \cite{RatraPeebles1988,Wetterich1988,Frieman1995}. Standard
quintessence potentials --- inverse power laws, exponentials,
pseudo-Nambu-Goldstone bosons --- are motivated by high-energy physics
(supersymmetry breaking, axion physics, string moduli) and share the
generic property that the field rolls toward infinity or has no
well-defined late-time attractor distinct from zero.

We study a specific quintessence potential:
\begin{equation}\label{eq:Vxi}
  V(\Xi) = \Lambda^{4}\,\Xi\log\Xi,
  \qquad \Xi > 0,\ \Lambda > 0,
\end{equation}
where $\Xi$ is a real, positive, dimensionless scalar field and $\Lambda$ is
a mass scale that sets the potential energy density. We shall call this the
\emph{logarithmic quintessence potential}. It has an analytic minimum at
$\Xi_0 = e^{-1}$. Because $\Xi$ is dimensionless and $\Lambda^{4}$ enters
as an overall prefactor, the location of this minimum is independent of
$\Lambda$; this is generic for any dimensionless-field potential of the
form $\Lambda^{4} f(\Xi)$ and is not a feature peculiar to $\Xi\log\Xi$.
The feature peculiar to this choice is the value $\Xi_0 = e^{-1}$, which
is also the maximizer of the per-bin Gibbs integrand $-x\log x$ over
$x \in (0,1]$. We discuss this structural connection in
\S\ref{sec:entropy} but do not treat it as derivational: $\Xi$ is an
unbounded real scalar field rather than a normalized probability density,
so the connection to thermodynamic entropy is a structural rhyme rather
than an operational identification. The model predicts freezing
quintessence dynamics; on the physical rolling branch studied numerically
the equation of state approaches $-1$ from above.

The paper is organized as follows. In \S\ref{sec:action} we state the
canonical action, fix unit conventions, and delimit the physical phase space.
In \S\ref{sec:derivations} we derive the Euler-Lagrange equation, the
stress-energy tensor, the analytic vacuum $\Xi_0 = e^{-1}$, the fluctuation
mass $m_\Xi^{2} = \Lambda^{4} e / M_\Pl^{2}$, and the bounded-below
global behavior of $V$.
Section \S\ref{sec:frw} specializes to a spatially flat FRW background and
derives the field energy density, pressure, and equation-of-state function.
Section \S\ref{sec:entropy} discusses the structural connection to the
Gibbs functional and to the Bialynicki-Birula--Mycielski nonlinearity in
nonlinear quantum mechanics. Section \S\ref{sec:obs} presents the
observational program, including an illustrative consistency check against
the DESI~2024 DR1 $(w_0, w_a)$ Gaussian summary, an explicit acknowledgment
of the methodological scope of that check, and discussions of EFT validity
over the field excursion and of the model's perturbation properties.
Section \S\ref{sec:comparison} compares the model to standard quintessence
potentials and to logarithmic prior art in the dark-sector literature.
Section \S\ref{sec:scope} states what the theory does \emph{not} claim.

\section{Canonical Action}\label{sec:action}

\subsection{Field content and conventions}\label{sec:fieldcontent}

The dynamical field content is
\[
  \bigl\{ g_{\mu\nu},\ \text{SM fields},\ \Xi(x) \bigr\},
\]
where $g_{\mu\nu}$ is a Lorentzian metric in the $(-,+,+,+)$ signature
convention, the SM sector follows the usual Standard Model prescription,
and $\Xi(x)$ is a real, positive, dimensionless scalar. $\Xi$ is a gauge
singlet and carries no charge under any internal symmetry, so the
minimal action contains no associated conserved current.

The effective theory is defined on the field domain $\Xi > 0$. In the
numerical analysis of \S\ref{sec:desifit} we restrict attention to initial
data $(\Xi_i, \dot\Xi_i)$ at $z_i \approx 20$ whose FRW evolution remains
within this domain over the redshift interval $0 \leq z \leq z_i$;
configurations leading to $\Xi \to 0^{+}$ during cosmic evolution are
outside the field domain considered in this effective theory and are
not considered.\footnote{An
equivalent formulation reparametrizes via $\psi = \log\Xi \in \mathbb{R}$,
eliminating the positivity constraint at the cost of a $\Xi$-dependent
kinetic prefactor $\propto e^{2\psi}$. A canonical dimension-1 field
$\phi = M_\Pl \Xi$ can also be introduced; in those variables the kinetic
term takes the standard form $\tfrac{1}{2}(\partial\phi)^{2}$ and the
potential becomes $V(\phi) = \Lambda^{4}(\phi/M_\Pl) \log(\phi/M_\Pl)$.
We work with the dimensionless field $\Xi$ throughout because in these
variables the stationary point $\Xi_0 = e^{-1}$ is especially transparent.}

\subsection{Action}

The minimal action is
\begin{equation}\label{eq:action}
  S = \int d^{4}x\,\sqrt{-g}\left[\frac{R}{16\pi G}
     + \mathcal{L}_{\SM}
     - \frac{M_\Pl^{2}}{2}\,g^{\mu\nu}\partial_\mu\Xi\,\partial_\nu\Xi
     - \Lambda^{4}\,\Xi\log\Xi\right],
\end{equation}
where $M_\Pl^{2} = (8\pi G)^{-1}$ is the reduced Planck mass and
$\Lambda$ is a mass scale that sets the potential energy density. The
signs of the kinetic and potential terms in \eqref{eq:action} are those
required by the $(-,+,+,+)$ signature convention adopted in
\S\ref{sec:fieldcontent}: for a homogeneous configuration $\Xi(t)$,
$g^{\mu\nu}\partial_\mu\Xi\,\partial_\nu\Xi = -\dot\Xi^{2}$, and the
displayed bracket evaluates to $\tfrac{1}{2}M_\Pl^{2}\dot\Xi^{2} -
V(\Xi)$, which is the standard non-ghost canonical form. Both
$\Xi$-dependent terms in the bracket have mass dimension 4: the kinetic
term acquires dimension $[M^{4}]$ from the $M_\Pl^{2}$ prefactor and the
two derivatives acting on the dimensionless $\Xi$, and the potential
term acquires dimension $[M^{4}]$ from the $\Lambda^{4}$ prefactor
multiplying the dimensionless combination $\Xi\log\Xi$. The single
model-specific parameter is $\Lambda$ (or equivalently the dimensionless
ratio $\Lambda^{4}/M_\Pl^{4}$); $\Lambda$ is fixed phenomenologically by
demanding the model reproduce the observed dark-energy density today,
which (as we show in \S\ref{sec:derivations}) places $\Lambda$
numerically near the dark-energy scale.

\begin{remark}
The model derives its predictions from $\Xi$ and $\Lambda$ alone; any
information-theoretic or entropic origin story for the potential form is
supplementary to the action and is deferred to \S\ref{sec:entropy}.
\end{remark}

\section{Derivations}\label{sec:derivations}

\subsection{Euler-Lagrange equation}

From $\mathcal{L}_\Xi = -\tfrac{1}{2} M_\Pl^{2}\, g^{\mu\nu}\partial_\mu\Xi\,\partial_\nu\Xi
- \Lambda^{4}\,\Xi\log\Xi$,
\[
  \frac{\partial \mathcal{L}_\Xi}{\partial \Xi} = -\Lambda^{4}(1 + \log\Xi),
  \qquad
  \nabla_\mu \frac{\partial \mathcal{L}_\Xi}{\partial(\nabla_\mu \Xi)}
  = -M_\Pl^{2}\,\Box\Xi.
\]
The Euler-Lagrange equation
$\partial\mathcal{L}_\Xi/\partial\Xi = \nabla_\mu(\partial\mathcal{L}_\Xi/\partial(\nabla_\mu\Xi))$
gives
\begin{equation}\label{eq:EL}
  \boxed{\;M_\Pl^{2}\,\Box\Xi = \Lambda^{4}(1 + \log\Xi)\;}
\end{equation}
or equivalently $\Box\Xi = (\Lambda^{4}/M_\Pl^{2})(1+\log\Xi)$, which
exhibits the characteristic mass scale $\Lambda^{4}/M_\Pl^{2}$ that
governs the dynamics. The vacuum condition $\Box\Xi = 0$ gives
$1 + \log\Xi_0 = 0$ as a coupling-independent algebraic equation, with
the prefactor $\Lambda^{4}/M_\Pl^{2}$ canceling at the vacuum.

\subsection{Stress-energy tensor}

Varying \eqref{eq:action} with respect to $g^{\mu\nu}$ gives
\begin{equation}\label{eq:stress}
  T^{\Xi}_{\mu\nu} = M_\Pl^{2}\,\partial_\mu\Xi\,\partial_\nu\Xi
  - g_{\mu\nu}\!\left[\frac{M_\Pl^{2}}{2}\,g^{\alpha\beta}\partial_\alpha\Xi\,\partial_\beta\Xi
  + \Lambda^{4}\,\Xi\log\Xi\right],
\end{equation}
manifestly symmetric in $(\mu, \nu)$. Covariant conservation
$\nabla^{\mu} T^{\Xi}_{\mu\nu} = 0$ holds on shell as a direct consequence
of \eqref{eq:EL} together with the contracted Bianchi identity, so
\eqref{eq:stress} is a consistent source for the Einstein equation
$G_{\mu\nu} = 8\pi G\, T^{\SM}_{\mu\nu} + 8\pi G\, T^{\Xi}_{\mu\nu}$.

\subsection{Vacuum/ground state}\label{sec:vacrenorm}

The potential $V(\Xi) = \Lambda^{4}\,\Xi\log\Xi$ has
\[
  V'(\Xi) = \Lambda^{4}(1 + \log\Xi),
  \qquad
  V''(\Xi) = \Lambda^{4}/\Xi,
\]
so $V'(\Xi_0) = 0$ at
\begin{equation}\label{eq:vacuum}
  \boxed{\;\Xi_0 = e^{-1} \approx 0.36788\;}
\end{equation}
and $V''(\Xi_0) = \Lambda^{4} e > 0$: $\Xi_0$ is a local minimum. The
location of this minimum at a coupling-independent dimensionless number
is generic for any potential of the form $\Lambda^{4} f(\Xi)$ with $\Xi$
dimensionless, since $V'(\Xi) = \Lambda^{4} f'(\Xi)$ vanishes
independently of $\Lambda$. What is specific to the choice
$f(\Xi) = \Xi\log\Xi$ is the value $\Xi_0 = e^{-1}$.

The potential energy at the vacuum is
\[
  V(\Xi_0) = \Lambda^{4}\,e^{-1} \log(e^{-1}) = -\Lambda^{4}/e < 0.
\]
This is the bare contribution from the scalar sector. As in standard
quintessence models, this combines with the bare cosmological constant
$\Lambda_{bare}$ from the gravitational sector to give the renormalized
observed cosmological constant $\Lambda_{obs} = \Lambda_{bare} - \Lambda^{4}/e$.
The bare quantities are not separately observable; only their sum enters
the Einstein equations. The scalar sector provides a dynamical
dark-energy component with a logarithmic potential and a vacuum
contribution whose observable effect is absorbed into the renormalized
cosmological constant. The model's empirical predictions for $w(z)$
depend on the dynamical evolution of $\Xi$ from $\Xi_i$ at $z_i$ toward
$\Xi_0 = e^{-1}$, not on the absolute value of $\Lambda_{bare}$.

\subsection{Fluctuation mass and stability}

Write $\Xi = \Xi_0 + \delta\Xi$ and linearize \eqref{eq:EL} about $\Xi_0$:
\[
  M_\Pl^{2}\,\Box\,\delta\Xi
  = \Lambda^{4}\,\frac{1}{\Xi_0}\,\delta\Xi
  = \Lambda^{4}\, e\,\delta\Xi.
\]
On a Minkowski background, a Fourier mode $\delta\Xi \propto e^{i(\mathbf{k}\cdot\mathbf{x} - \omega t)}$
satisfies $\omega^{2} - k^{2} = \Lambda^{4} e / M_\Pl^{2}$, yielding the
Klein-Gordon dispersion $\omega^{2} = k^{2} + m_\Xi^{2}$ with
\begin{equation}\label{eq:mass}
  \boxed{\;m_\Xi^{2} = \frac{\Lambda^{4}\,e}{M_\Pl^{2}}\;}
\end{equation}
Stability: $m_\Xi^{2} > 0$ for any $\Lambda > 0$.

\paragraph{Dark-energy scale.}
For $\Xi$ to play the role of dark energy with mass of order the Hubble
scale today, $m_\Xi \sim H_0 \sim 10^{-33}$ eV, equation \eqref{eq:mass}
implies $\Lambda \sim (H_0\, M_\Pl/e^{1/2})^{1/2}$, which evaluates to
$\Lambda \approx 1.5$~meV from this dimensional analysis. The numerical
fit reported in \S\ref{sec:desifit} gives $\Lambda \approx 1.7$~meV. The
ratio of the two values is set by the geometric factor relating
$m_\Xi^{2} = \Lambda^{4}e/M_\Pl^{2}$ to $\rho_{c,0} = 3 H_0^{2}M_\Pl^{2}$:
$\Lambda_\mathrm{fit}/\Lambda_\mathrm{dim} = (3 e \Omega_\Lambda)^{1/4}
\approx 1.13$, in agreement with the empirical ratio $1.7/1.5 \approx 1.13$.
Both values place $\Lambda$ near the observed dark-energy scale; this is
a consistency observation, not a derivation.

\subsection{Global behavior of $V$}

For $\Xi > 0$, $V(\Xi) = \Lambda^{4}\,\Xi\log\Xi$ satisfies
\[
  V(\Xi) \to 0\ \text{from below as}\ \Xi \to 0^{+},\qquad
  V(\Xi_0) = -\Lambda^{4}/e\ (\text{global minimum}),\qquad
  V(\Xi) \to +\infty\ \text{as}\ \Xi \to +\infty.
\]
The scalar potential is globally bounded below on $\Xi > 0$, and the
homogeneous scalar-sector energy density
$\rho_\Xi = \tfrac{M_\Pl^{2}}{2}\dot\Xi^{2} + V(\Xi)$ is bounded below
by $-\Lambda^{4}/e$.

\section{FRW Cosmology}\label{sec:frw}

\subsection{Background equations}

In a spatially flat FRW background with metric
$ds^{2} = -dt^{2} + a^{2}(t)\, d\mathbf{x}^{2}$ and homogeneous $\Xi(t)$,
the kinetic invariant is
$g^{\alpha\beta}\partial_\alpha \Xi\, \partial_\beta \Xi = -\dot\Xi^{2}$.
Writing $T^{\Xi}_{\mu\nu} = \diag(\rho_\Xi, a^{2} p_\Xi, a^{2} p_\Xi, a^{2} p_\Xi)$,
\begin{align}
  \rho_\Xi &= \frac{M_\Pl^{2}}{2}\,\dot\Xi^{2} + \Lambda^{4}\,\Xi\log\Xi,\label{eq:rho}\\
  p_\Xi &= \frac{M_\Pl^{2}}{2}\,\dot\Xi^{2} - \Lambda^{4}\,\Xi\log\Xi.\label{eq:p}
\end{align}
The Friedmann equations are
\begin{equation}\label{eq:friedmann}
  H^{2} = \frac{8\pi G}{3}\!\left(\rho_{\SM} + \rho_\Xi\right),\qquad
  \dot H = -4\pi G\!\left(\rho_{\SM} + p_{\SM} + M_\Pl^{2}\,\dot\Xi^{2}\right),
\end{equation}
and the $\Xi$ equation of motion, from \eqref{eq:EL} in FRW (where
$\Box\Xi = -\ddot\Xi - 3H\dot\Xi$ in the $(-,+,+,+)$ convention), is
\begin{equation}\label{eq:xieom}
  \boxed{\;M_\Pl^{2}(\ddot\Xi + 3H\dot\Xi) = -\Lambda^{4}(1 + \log\Xi)\;}
\end{equation}

\subsection{Equation of state}

From \eqref{eq:rho}-\eqref{eq:p},
\begin{equation}\label{eq:w}
  w_\Xi = \frac{p_\Xi}{\rho_\Xi}
  = \frac{\tfrac{1}{2} M_\Pl^{2}\dot\Xi^{2} - \Lambda^{4}\,\Xi\log\Xi}
         {\tfrac{1}{2} M_\Pl^{2}\dot\Xi^{2} + \Lambda^{4}\,\Xi\log\Xi}
  = -1 + \frac{M_\Pl^{2}\dot\Xi^{2}}{\rho_\Xi}.
\end{equation}

\begin{proposition}[Limiting regimes]\label{prop:limits}
For $\Xi > 0$ the EoS \eqref{eq:w} satisfies
\begin{enumerate}
  \item[(i)] $w_\Xi = -1$ at the vacuum ($\dot\Xi = 0$, $\Xi = \Xi_0$):
        vacuum-like equation of state.
  \item[(ii)] On the physical rolling branch studied numerically
        (\S\ref{sec:desifit}), $w_\Xi$ approaches $-1$ from above as
        $\Xi \to \Xi_0$ and $\dot\Xi \to 0$.
  \item[(iii)] $w_\Xi \to +1$ in the kinetic-domination limit
        ($M_\Pl^{2}\dot\Xi^{2} \gg \Lambda^{4}|\Xi\log\Xi|$): stiff fluid.
\end{enumerate}
\end{proposition}

\begin{proof}
(i) Direct substitution in \eqref{eq:w}: $p_\Xi = -\Lambda^{4}/e$ and
$\rho_\Xi = \Lambda^{4}\,\Xi_0 \log\Xi_0 = -\Lambda^{4}/e$, giving
$w_\Xi = -1$. (ii) On the monotone rolling solutions studied numerically,
with $\Xi > \Xi_0$ and $\rho_\Xi > 0$ over the fitted redshift range, the
kinetic term decreases as the field approaches the minimum, and
\eqref{eq:w} gives $w_\Xi \to -1$ from above.
(iii) In the kinetic-domination regime, both numerator and denominator
of \eqref{eq:w} are dominated by $M_\Pl^{2}\dot\Xi^{2}$, giving
$w_\Xi \to +1$.
\end{proof}

\begin{remark}[Bare vs renormalized vacuum]
The equality $w_\Xi = -1$ at $\Xi = \Xi_0$ in (i) is a formal statement
about the scalar-sector stress-energy tensor before vacuum-energy
renormalization; the physical late-time vacuum entering cosmology is the
renormalized sum discussed in \S\ref{sec:vacrenorm}.
\end{remark}

\begin{remark}[On phantom crossing]
A general structural theorem $w_\Xi \geq -1$ would require careful
treatment of the bare scalar-sector energy density, which is negative
at the vacuum ($\rho_\Xi(\Xi_0) = -\Lambda^4/e$). Over the fitted
rolling branch in \S\ref{sec:desifit}, the numerical solution exhibits
$w_\Xi(z) \geq -1$ at all redshifts and produces no phantom crossing;
we treat the absence of phantom crossing as a property of the fitted
branch rather than a global theorem of the minimal theory.
\end{remark}

\subsection{Late-time attractor}

Combining \eqref{eq:xieom} with \eqref{eq:friedmann}, the field rolls
toward the vacuum, and the damping term $3H\dot\Xi$ drives $\dot\Xi \to 0$
as $H$ remains positive. In this limit $\Xi \to \Xi_0 = e^{-1}$ and, by
Proposition \ref{prop:limits}(i), $w_\Xi \to -1$. The
late-time attractor is therefore vacuum-like behavior with $w \to -1$.

\section{Structural Connections to Information-Theoretic Functionals}\label{sec:entropy}

The functional form $\Xi\log\Xi$ that appears in the potential
\eqref{eq:Vxi} is recognizable from two adjacent settings: the per-bin
integrand of the Gibbs/Shannon entropy on a normalized probability
distribution, and the Bialynicki-Birula--Mycielski nonlinearity in
nonlinear quantum mechanics. We treat both connections as structural
observations about the functional form of $V$, not as derivations of
\eqref{eq:action}, and we are explicit below about where the analogies
break down.

\subsection{The Gibbs functional}

The function
\begin{equation}\label{eq:entropy}
  H_{\Gibbs}(x) = -x \log x,
  \qquad x \in (0,1],
\end{equation}
is the per-bin contribution to the Gibbs/Shannon entropy of a discrete
probability distribution $\{p_i\}$ via $H = \sum_i H_{\Gibbs}(p_i)$,
where the $p_i$ are nonnegative and sum to $1$. The function attains its
maximum on $(0,1]$ at $x = e^{-1}$, with $H'_{\Gibbs}(e^{-1}) = 0$ and
$H''_{\Gibbs}(e^{-1}) = -e < 0$.

The potential \eqref{eq:Vxi} has the same functional form up to a sign
and the dimensionful prefactor:
\begin{equation}\label{eq:formal-rel}
  V(\Xi) = -\Lambda^{4}\,H_{\Gibbs}(\Xi).
\end{equation}
The vacuum condition $V'(\Xi_0) = 0$ is therefore equivalent to
$H'_{\Gibbs}(\Xi_0) = 0$, and the same $\Xi_0 = e^{-1}$ that minimizes
$V$ also maximizes $H_{\Gibbs}$ on $(0,1]$.

\begin{remark}[Where the analogy breaks]
The function $H_{\Gibbs}$ as written has an operational interpretation
as a thermodynamic entropy only when $x$ is a normalized probability
($x \in [0,1]$ with $\sum x = 1$). In the present scalar-field theory
$\Xi$ is a real, positive, unbounded scalar with no normalization
constraint. The relation \eqref{eq:formal-rel} is therefore a
structural rhyme between $V$ and the Gibbs integrand, not an
identification of $-V/\Lambda^{4}$ with a thermodynamic entropy.
We do not claim that the field's dynamics is equivalent to entropy
maximization in any operational sense; we record only that the
potential's functional form coincides with the Gibbs integrand and
that $\Xi_0 = e^{-1}$ coincides numerically with the per-bin
entropy maximum. The model's empirical predictions for $w(z)$
\S\ref{sec:obs} do not depend on this connection.
\end{remark}

\subsection{Bialynicki-Birula--Mycielski nonlinearity}

The same logarithmic functional form appears in nonlinear quantum
mechanics: Bialynicki-Birula and Mycielski
\cite{BialynickiBirulaMycielski1976} characterized
$|\psi|^{2} \log |\psi|^{2}$ as the unique analytic nonlinear extension
of the free Schr\"odinger equation that preserves separability of the
density $|\psi|^{2}$ under wave-function factorization
$\psi = \psi_1 \psi_2$.

\begin{remark}[Where the analogy breaks]
The Bialynicki-Birula--Mycielski nonlinearity is in
$|\psi|^{2}\log|\psi|^{2}$, where $|\psi|^{2}$ is a probability density.
In the present setting, the analogous variable is the field $\Xi$
itself, not $\Xi^{2}$, and $\Xi$ has no probability-density
interpretation. Importing the same functional form from a setting
where the argument is a normalized squared amplitude to a setting where
the argument is an unbounded real scalar is a structural rhyme, not a
derivation of \eqref{eq:action} from the BBM uniqueness theorem. We
cite \cite{BialynickiBirulaMycielski1976,Rosen1969,CazenaveHaraux1980,HoeghKrohn1971}
as motivation for studying $V \propto \Xi\log\Xi$ as a candidate
quintessence potential, not as a uniqueness proof for it.
\end{remark}

\subsection{Information-theoretic precedents}

\begin{remark}[Information-theoretic dark energy]
Entropic approaches to dark energy in the holographic sense
\cite{Verlinde2011,CHLZ2012} typically start from an entropy ansatz
and derive a dark-energy density. The present approach is structurally
different: the action \eqref{eq:action} is specified first, and the
relation \eqref{eq:formal-rel} between $V$ and the Gibbs integrand
follows. Information-theoretic field theory \cite{Ensslin2013,Caticha2012}
uses $\Xi\log\Xi$ as a standard KL-divergence integrand on spaces of
probability densities; the appearance of the same functional as a
cosmological self-interaction is a point of contact between these
literatures and the present work, but the connection is again
structural rather than derivational.
\end{remark}

The model's empirical content (\S\ref{sec:obs}) does not depend on either
of the two structural connections discussed in this section.

\section{Observational Predictions}\label{sec:obs}

\subsection{The $w(z)$ trajectory}

On the physical rolling branch studied numerically, we write
$w_\Xi(z) = -1 + \epsilon(z)$ with $\epsilon(z) \geq 0$. The freezing
attractor derived above gives $\epsilon(z) \to 0$ at late times, with
the rate determined by $\Lambda$ and the initial displacement $\Xi(z_i)$
at some reference redshift $z_i$. DESI DR1/DR2 measurements of $w_0$
and $w_a$ in the CPL parametrization probe $\epsilon(z)$ directly at
$z \lesssim 1.5$ \cite{DESI2024BAO,DESI2024VI}.

\subsection{Numerical integration and DESI consistency check}\label{sec:desifit}

We perform an illustrative background-level consistency check against
the DESI~2024 DR1 published $(w_0, w_a)$ central values, treating those
values as a Gaussian summary on the CPL parametrization. We emphasize at
the outset what this check does and does not do: the
$\chi^{2}$ values reported below are computed against the published
$(w_0, w_a)$ marginal Gaussian summary, with central values
$w_0 = -0.827 \pm 0.063$, $w_a = -0.75 \pm 0.27$ \cite{DESI2024VI}, and
they quantify how close the model's $(w_0, w_a)$ output sits to those
published central values \emph{under that Gaussian-on-summary
approximation}. They are not derived from the underlying joint
BAO + CMB + SN likelihood. The DR1 BAO data alone are consistent with
flat $\Lambda$CDM \cite{DESI2024BAO}; the dynamical-dark-energy
preference at the $2.8$--$4.2\sigma$ level appears in the joint
likelihood, not in the marginal Gaussian on $(w_0, w_a)$. A full joint
likelihood analysis using the DR2 chains \cite{DESI2025DR2} together
with Pantheon+ supernovae and Planck CMB constraints is in preparation
as a companion numerical paper.

We numerically integrate the FRW system
\eqref{eq:friedmann}-\eqref{eq:xieom} over $z \in [0, 10]$ with a
standard-density background ($\Omega_m = 0.315$, $\Omega_r = 9.1\times10^{-5}$,
$\Omega_\Xi = 0.6849$, $H_0 = 67.4$ km s$^{-1}$ Mpc$^{-1}$
\cite{Planck2020}) and scan over $(\Lambda^{4}/\rho_{c,0}, \Xi_i, \dot\Xi_i)$
at initial redshift $z_i \approx 20$, where $\rho_{c,0} = 3 H_0^{2} M_\Pl^{2}$
is the critical density today.

\paragraph{Energy-budget framing.}
The reported $\Omega_\Xi = 0.6849$ is the dark-energy budget of the
combined sector --- renormalized cosmological constant
$\Lambda_\mathrm{obs} = \Lambda_\mathrm{bare} - \Lambda^{4}/e$ from
\S\ref{sec:vacrenorm} together with the scalar's dynamical
deviation $V(\Xi) - V(\Xi_0)$ and its kinetic energy --- and the
reported $(w_0, w_a)$ describe the effective equation of state of
this combined sector against DR1. With
$\Xi(z=0) \approx 0.97$ above $\Xi_0 = e^{-1}$ and
$\Lambda^{4}/\rho_{c,0} = 0.231$, the scalar's dynamical
deviation contributes $\Lambda^{4}\,(\Xi\log\Xi + 1/e)/\rho_{c,0}
\approx 0.078\,\rho_{c,0}$ in potential energy at the present
epoch, with the remainder $\Omega_{\Lambda_\mathrm{obs}}
\approx 0.61\,\rho_{c,0}$ carried by the renormalized constant.
The script $\texttt{desi\_xi\_optimize\_v2.py}$ tracks the scalar
field's intrinsic $\rho_\Xi$ and $w_\Xi$ as defined in
\eqref{eq:rhoxi}-\eqref{eq:wxi}; the comparison to DR1 in
the table below should be read as a proximity check on the
combined-sector effective $w(z)$ under the renormalization
described in \S\ref{sec:vacrenorm}, not on the scalar field's
intrinsic equation of state, which is dominated by the kinetic
component on this configuration. A full joint
likelihood that separates $\Lambda_\mathrm{obs}$ from the scalar
sector self-consistently is deferred to the companion numerical
paper. We restrict to the freezing branch of the
model, implemented by requiring a monotone $w(z)$ trajectory over the
fitted range (criterion (F2) of \S\ref{sec:falsify}). A configuration
that yields proximity to the DESI DR1 published $(w_0, w_a)$ Gaussian
summary is
\begin{equation}\label{eq:desibestfit}
  \Lambda^{4}/\rho_{c,0} = 0.231,
  \quad \Xi_i = 0.925,
  \quad \dot\Xi_i = +0.429,
\end{equation}
corresponding to $\Lambda \approx 1.7\,\text{meV}$ in physical units and
yielding $w_0^{\Xi} = -0.793$ and $w_a^{\Xi} = -0.451$. Two
$\chi^{2}$ values quantify the proximity of these output values to
the published DESI DR1 $(w_0, w_a)$ marginal Gaussian summary:
\[
  \chi^{2}_{\text{Gauss},\rho=0} = 1.52,\qquad
  \chi^{2}_{\text{Gauss},\rho=-0.9} = 13.6,
\]
both computed against the published central values $w_0 = -0.827 \pm
0.063$, $w_a = -0.75 \pm 0.27$. The first ignores the $(w_0, w_a)$
correlation present in the DR1 posterior; the second uses the
published correlation $\rho(w_0, w_a) \approx -0.9$ via
$\chi^{2}_{\rho} = (1-\rho^{2})^{-1}\bigl(x^{2}-2\rho x y + y^{2}\bigr)$
with $x=(w_0-w_0^{\text{DR1}})/\sigma_{w_0}$ and analogously for $y$.
For comparison, $\Lambda$CDM at $(w_0, w_a) = (-1, 0)$ gives
$\chi^{2}_{\text{Gauss},\rho=0,\Lambda\text{CDM}} = 15.26$ and
$\chi^{2}_{\text{Gauss},\rho=-0.9,\Lambda\text{CDM}} = 8.03$.
The two metrics give qualitatively different pictures: under the
uncorrelated approximation the $\Xi$ configuration above sits
$\sim 1.2\sigma$ from the DR1 central values and $\Lambda$CDM sits
$\sim 3.9\sigma$ away; under the correlated metric both the
$\Xi$ configuration and $\Lambda$CDM sit a few $\sigma$ from the
DR1 central, with $\Lambda$CDM marginally closer
($\Delta\chi^{2}_{\text{Gauss},\rho=-0.9}
= \chi^{2}_{\Xi} - \chi^{2}_{\Lambda\text{CDM}} \approx 5.6$ on
2 d.o.f., a $\sim 2.4\sigma$ disfavoring of the documented
$\Xi$ configuration relative to $\Lambda$CDM under the
Gaussian-on-summary approximation). The correlated
metric is the disciplined comparison; the uncorrelated $\chi^{2}$ is
reported only because it appeared in earlier drafts of this
manuscript and may persist as a default in similar consistency
checks elsewhere in the literature. Neither metric is a substitute
for the joint $\text{BAO} + \text{CMB} + \text{SN}$ likelihood,
which is deferred to the companion numerical paper. The
$\Xi$-model's $w(z)$ trajectory on this configuration is monotone
freezing:

\begin{center}
\begin{tabular}{c|c}
$z$ & $w_\Xi(z)$\\
\hline
$0.0$ & $-0.7933$\\
$0.3$ & $-0.9098$\\
$0.5$ & $-0.9432$\\
$0.8$ & $-0.9701$\\
$1.0$ & $-0.9804$\\
$1.5$ & $-0.9946$\\
$2.0$ & $-0.9998$
\end{tabular}
\end{center}

The present-epoch field value is $\Xi(z=0) \approx 0.97$, still above
the vacuum $\Xi_0 = e^{-1} \approx 0.37$, confirming the field is in
its rolling phase. The numerical integration is performed by the
script \texttt{desi\_xi\_optimize\_v2.py} provided as electronic
supplementary material.

\paragraph{On the ``freezing'' label.}
We use the term \emph{freezing} in the asymptotic sense
($w(z)\to -1$ as the field eventually settles at $\Xi_0$). Over
the observed range $z\in[0, 2]$ the field is on an \emph{outbound}
branch with $\Xi$ rising from $\Xi_i = 0.925$ at $z\approx 20$ to
$\Xi(0)\approx 0.97$, with the turnaround and roll-back to vacuum
lying in the future. The equation of state $w_\Xi(z)$ in the
table above is monotone in $z$ on this range, which is the
``freezing-branch'' criterion (F2) of \S\ref{sec:falsify} that we
use to select solutions; it is consistent with the standard
freezing-quintessence terminology in the asymptotic limit but
does not require the field to be inbound during the
observable window.

\paragraph{Initial-condition tuning.}
The values $(\Xi_i, \dot\Xi_i) = (0.925, +0.429)$ at $z_i \approx 20$ in
\eqref{eq:desibestfit} are tuned, not predicted. They are selected by
the requirement of proximity to the DESI DR1 $(w_0, w_a)$ central values
together with the freezing-branch criterion $w(z)$ monotone. We note
explicitly that $\dot\Xi_i > 0$ corresponds to an \emph{outbound}
initial condition: at $z_i$ the field is moving away from the
vacuum $\Xi_0 = e^{-1}$, decelerates under Hubble friction and
potential gradient, and rolls back toward $\Xi_0$ at later times. The
direction of initial motion is itself a tuning choice in addition to
the values $(\Xi_i, \dot\Xi_i)$. There is
no naturalness or attractor argument in the present minimal theory
that singles out these values; the freezing-branch restriction selects
solutions, it does not derive them. This sensitivity to initial
conditions is shared with most quintessence models that fit
late-time data, and a complete naturalness analysis (e.g., a
quintessence-tracker construction) is outside the scope of the present
paper. A reduction of this two-dimensional initial-condition freedom
to a one-dimensional or zero-dimensional family --- e.g., via a
late-time attractor mechanism that erases initial-condition memory ---
would strengthen the model and is a target for follow-up work.

\paragraph{Methodological scope of this consistency check.}
To restate the scope explicitly: the $\chi^{2}_{\text{Gauss}}$ values
reported here are computed against published CPL summary statistics,
treated as an uncorrelated Gaussian likelihood on $(w_0, w_a)$. The
DESI DR1/DR2 publications provide non-negligible $(w_0, w_a)$
correlation in their full chains; ignoring that correlation is one
limitation of the present check. A full joint likelihood analysis
on the underlying BAO + CMB + SN data is deferred to a companion
numerical paper. The present check demonstrates that the model can
produce $(w_0, w_a)$ values within $\sim 1\sigma$ of the DESI DR1
central values for $\Lambda \approx 1.7$~meV with appropriately
chosen initial conditions, and that $\Lambda$CDM sits roughly
$3.9\sigma$ from those central values under the same Gaussian
approximation. It is not a claim of best-fit to the full DESI
likelihood, nor a measure of the model's preference relative to
$\Lambda$CDM in any joint-likelihood sense.

\paragraph{Reproducibility.}
The numerical fit reported above is reproduced by a single command. The
script \texttt{desi\_xi\_optimize\_v2.py} (electronic supplementary material)
accepts no arguments; running it produces the best-fit parameters
\eqref{eq:desibestfit}, the $w(z)$ trajectory, and the
$\chi^{2}_{\text{Gauss}}$ values for both the $\Xi$-model and
$\Lambda$CDM. Total runtime is under 30 seconds on a standard laptop.
The verification script \texttt{proof\_xi\_canonical.py} runs 22
algebraic and stability tests on the field equations; all pass. Both
scripts are archived at DOI~\texttt{10.5281/zenodo.18852047}.

\subsection{Falsification criteria}\label{sec:falsify}

The model predicts:
\begin{enumerate}
  \item[(F1)] On the monotone rolling branch selected by the present
        model, $w_\Xi(z) \geq -1$ over the redshift range studied
        numerically. A robust observational preference for $w(z) < -1$
        across that same range, at $\geq 3\sigma$ confidence in the joint
        DR2 + Pantheon+ + Planck likelihood, would disfavor this branch
        of the model.
  \item[(F2)] Monotone approach $w(z) \to -1$ from above (freezing). A
        confirmed non-monotone $w(z)$ trajectory over $0 \leq z \leq 1.5$
        at $\geq 3\sigma$ in the joint DR2 + Pantheon+ + Planck likelihood
        would falsify the rolling-branch fit.
  \item[(F3)] Asymptotic vacuum endpoint $w \to -1$. The present-epoch
        value $w_0 \neq -1$ in general while the field is still rolling;
        the manuscript's best fit gives $w_0 \approx -0.79$ at the present
        epoch. Evidence favoring late-time behavior inconsistent with
        $w \to -1$, at $\geq 3\sigma$ in the joint Stage-IV likelihood,
        would disfavor the freezing attractor interpretation.
  \item[(F4)] Specific $w(z)$ profile fixed by two parameters
        $(\Lambda^{4}, \Xi_i)$. The documented fit predicts
        $w_\Xi(z = 0.5) = -0.943 \pm 0.02$, where the
        $\pm 0.02$ tolerance reflects the spread of the fit over the
        documented 1$\sigma$ region of $(\Lambda^{4}/\rho_{c,0}, \Xi_i)$
        in \texttt{desi\_xi\_optimize\_v2.py}. After fitting $w_0$ and
        $w_a$ at $z = 0$, the value of $w(z = 0.5)$ is overdetermined
        in the model's two-parameter family. A measured
        $w(z = 0.5)$ outside $-0.943 \pm 0.05$ at Stage-IV precision
        ($\sigma(w(z = 0.5)) \approx 0.03$ per
        \cite{ShajibFrieman2025}) would disfavor the documented fit.
\end{enumerate}
Stage-IV surveys (DESI through DR3, Euclid, Rubin LSST) will reduce
$\sigma(w_0)$ to approximately $0.02$ and $\sigma(w_a)$ to approximately
$0.10$ \cite{ShajibFrieman2025,Turyshev2026Review}, providing the
precision needed to test (F1)-(F4) at the levels indicated above.

\subsection{Fifth-force and local tests}

In \eqref{eq:action}, $\Xi$ has no direct matter coupling; any effective
coupling to SM matter arises only through gravity and is Planck-suppressed.
Such gravitational-strength couplings are well below the sensitivity of
laboratory tests of the inverse-square law \cite{EotWash,MICROSCOPE},
which probe direct-matter couplings several orders of magnitude above
the gravitational scale. The minimal theory is therefore expected to
produce \emph{no} observable fifth force and is consistent with all
current local tests. Any direct-matter coupling (conformal
$g\Xi\, T^{\SM}$, disformal $f(\Xi) T^{\SM}$, or universal rescaling
$(\Xi/\Xi_0) \mathcal{L}_{\SM}$) is outside the scope of the minimal
action and would be added as a separate extension.

\subsection{Effective-field-theory validity over the field
excursion}\label{sec:eft}

In canonical variables $\phi = M_\Pl \Xi$ (cf.\ the footnote in
\S\ref{sec:fieldcontent}), the vacuum sits at $\phi_0 = M_\Pl/e$ and
the fitted trajectory has $\phi$ traversing field-space distances of
order $M_\Pl$. This is the standard quintessence transplanckian-field
concern: an EFT obtained by integrating out heavier degrees of freedom
at a scale $M$ generically receives operator corrections suppressed by
powers of $\phi/M$, and for $\phi \sim M_\Pl$ such corrections are not
manifestly under control unless the underlying theory protects them
(via approximate shift symmetry, supersymmetry, or some other
mechanism). The minimal action \eqref{eq:action} does not exhibit a
shift symmetry, since $\Xi\log\Xi$ is not invariant under
$\Xi \to \Xi + c$, and we do not propose an underlying UV completion in
the present paper.

We make three observations. (i) The field excursion in canonical
variables is transplanckian in absolute terms: at the documented fit
$(\Xi_i, \dot\Xi_i) = (0.925, +0.429)$ at $z_i \approx 20$, the
field traverses $\Delta\phi = M_\Pl(\Xi_i - \Xi_0) \approx 0.56\,M_\Pl$
between the initial epoch and the present, with the late-time
($0 \le z \le 2$) portion of the trajectory contributing
$\Delta\phi \approx 0.20\,M_\Pl$. The argument that the
EFT remains controlled cannot rely on a small absolute field
excursion. We do not have a UV completion that protects the form
of $V$ against $M_\Pl$-scale operator corrections, and the minimal
action \eqref{eq:action} does not exhibit a shift symmetry that
would suppress them in the absence of such a completion. The
present paper treats \eqref{eq:action} as a tree-level effective
description; the question of UV control over the transplanckian
field excursion is open. This is the standard transplanckian
issue for any quintessence model with $\Delta\phi \sim M_\Pl$ and
is shared with most freezing-quintessence constructions in the
literature.
(ii) The mass scale $\Lambda \approx 1.7$~meV is many
orders of magnitude below $M_\Pl$, and the dimensional analysis
$m_\Xi^{2} = \Lambda^{4} e / M_\Pl^{2} \sim H_0^{2}$ does not require
new physics at scales below $M_\Pl$. (iii) A controlled UV embedding
of \eqref{eq:action} is a standard open question for any tree-level
quintessence potential and is not addressed here. We flag this as a
limitation and leave it to subsequent work; the present paper studies
the phenomenology of \eqref{eq:action} treated as a low-energy
effective description.

\subsection{Perturbations}\label{sec:perts}

The minimally coupled canonical scalar field \eqref{eq:action} has
\emph{rest-frame} sound speed
\[
  c_{s}^{2} \;=\; 1
\]
for the propagation of the perturbation $\delta\Xi$, by the standard
result for canonical scalars. This is
distinct from the \emph{adiabatic} sound speed
$c_{a}^{2} = \dot p_\Xi/\dot\rho_\Xi$, which depends on $w(z)$ and is
generally not equal to unity for a rolling quintessence field. The
quantity controlling clustering on subhorizon scales is $c_{s}^{2}$,
not $c_{a}^{2}$, by direct evaluation of \eqref{eq:rho}-\eqref{eq:p}
in the rest frame of the perturbation.
The model therefore predicts the standard quintessence behavior: dark-energy
perturbations are sound-speed-suppressed on subhorizon scales, and the
field does not cluster appreciably on cosmological scales of present
observational interest. On scales $k \gg a H$ the perturbations $\delta\Xi$
oscillate with frequency $c_s k/a = k/a$ and are damped relative to matter
perturbations; the contribution of $\delta\Xi$ to the late-time matter
power spectrum at $k \sim 0.1$~$h$/Mpc is therefore at the percent level
and below current Stage-IV sensitivity for distinguishing $c_s^{2} = 1$
quintessence from $\Lambda$CDM at the perturbation level
\cite{ShajibFrieman2025,Turyshev2026Review}. A full perturbation
calculation including ISW and CMB lensing predictions, and any departures
from $c_s^{2} = 1$ that arise from non-minimal couplings, is deferred to
the companion numerical paper.

\subsection{Gravitational-wave speed}

In \eqref{eq:action}, tensor perturbations propagate at the speed of
light, in accord with GW170817 \cite{GW170817}. No superluminal or
birefringent signal is predicted.

\section{Comparison to Literature}\label{sec:comparison}

Standard quintessence potentials, their existence of a minimum, their
late-time EoS behavior, and their comparison to the present model are
summarized in Table~\ref{tab:compare}.

\begin{table}[h]
\small
\begin{tabular}{l|c|c|c|l}
Potential & Reference & Minimum? & $w \to$ at late time & Notes\\
\hline
$V \sim \phi^{-\alpha}$ & \cite{RatraPeebles1988} & no & $w^{*} > -1$ const & tracker\\
$V \sim e^{-\alpha\phi}$ & \cite{Wetterich1988} & no & $w^{*}$ const & rolling\\
$V \sim 1 + \cos(\phi/f)$ & \cite{Frieman1995} & yes ($\phi = \pi f$) & approximate $-1$ & PNGB / thawing\\
$V_\text{CW}(\phi)$ & \cite{ColemanWeinberg1973} & yes, $\langle\phi\rangle(\lambda)$ & --- & loop-level $\phi^{4}\log\phi$, distinct family\\
$V \sim \phi^{p}(\log\phi)^{q}$ & \cite{BarrowParsons1995} & depends on $(p,q)$ & --- & inflation family\\
$\boxed{V = \Lambda^{4}\Xi\log\Xi}$ & present work & yes ($\Xi_0 = e^{-1}$) & $-1$ asymptotically & freezing, $V = -\Lambda^{4} H_{\Gibbs}$
\end{tabular}
\caption{Comparison of quintessence potentials. The last row is the
model of the present paper. The closest structural relative is the
Barrow-Parsons inflation family \cite{BarrowParsons1995}
$V \sim V_0 \phi^{p}(\log\phi)^{q}$, which contains the present form as
the $(p,q) = (1,1)$ subcase but is applied to inflation, not dark energy,
and does not single out this subcase or derive its vacuum at $e^{-1}$.}
\label{tab:compare}
\end{table}

\begin{remark}[On novelty]
Logarithmic functional forms appear in the dark-sector literature in
several distinct contexts. Thompson \cite{Thompson2019} studied
logarithmic and exponential dark-energy potentials in a beta-function
quintessence reconstruction framework, with logarithmic potentials
arising from the beta-function formalism rather than as fundamental
quintessence inputs. Kouwn and Oh \cite{KouwnOh2012} proposed a
dark-energy fluid with logarithmic equation of state, with low-energy
effective potential $V(\phi) = \exp(-2\phi)\bigl(\tfrac{4A}{3}\phi + B\bigr)$.
In an adjacent direction, Chavanis \cite{Chavanis2015} developed a
logotropic dark-fluid program addressing dark-matter and dark-energy
unification, with logarithmic structure in the dark-sector equation of
state. Oikonomou \cite{Oikonomou2019} analyzed coupled dark-sector
systems with generalized logarithmic equations of state. The present
work is distinct from each of these in studying a specific dimensionless
scalar with quintessence potential $V = \Lambda^{4}\,\Xi\log\Xi$,
deriving an exact analytic vacuum at $\Xi_0 = e^{-1}$ that is independent
of the coupling, and obtaining freezing equation-of-state behavior with
asymptotic limit $w \to -1$. To the best of our
literature search
(covering arXiv categories astro-ph.CO, hep-ph, and gr-qc through April
2026, including DR2-era $V(\phi)$ reconstructions and analytic-family
surveys
\cite{AdamEtAl2025,Adil2026Analytic,Wang2026QReconstruction}), the
specific combination of features --- dimensionless field, the exact
$\Xi\log\Xi$ functional form, the coupling-independent $e^{-1}$ vacuum,
and the freezing equation-of-state behavior with exact asymptotic
endpoint $w \to -1$ --- has not been previously
studied as a quintessence model.
\end{remark}

\section{Scope and Limitations}\label{sec:scope}

The present paper asserts the following, and only the following:
\begin{itemize}
  \item The action \eqref{eq:action} is a minimal and self-consistent
        scalar-tensor theory. Its Euler-Lagrange equation \eqref{eq:EL},
        stress-energy tensor \eqref{eq:stress}, vacuum \eqref{eq:vacuum},
        fluctuation mass \eqref{eq:mass}, and FRW equations
        \eqref{eq:friedmann}-\eqref{eq:xieom} are formal consequences.
  \item The vacuum $\Xi_0 = e^{-1}$ coincides with the maximizer of the
        per-bin Gibbs integrand $-x\log x$ on $(0,1]$. As detailed in
        \S\ref{sec:entropy}, this is a structural relation between
        the functional form of $V$ and that of $H_{\Gibbs}$, not an
        identification of $\Xi$ with a probability density or of
        $-V/\Lambda^{4}$ with a thermodynamic entropy.
  \item The DESI consistency check reported in \S\ref{sec:desifit} is a
        numerical integration of \eqref{eq:friedmann}-\eqref{eq:xieom}
        against the published DESI DR1 (2024) $(w_0, w_a)$ central
        values; the best-fit parameters \eqref{eq:desibestfit} are not
        the result of a full likelihood analysis including BAO, CMB,
        and SN data jointly, which is deferred to a companion numerical
        paper.
\end{itemize}

The present paper does \emph{not} claim the following:
\begin{itemize}
  \item Any derivation of $\Lambda$ from a microphysical substrate; the
        observation that $\Lambda$ is numerically near the dark-energy
        scale when $m_\Xi \sim H_0$ is a consistency check, not a
        solution to the cosmological-constant hierarchy problem.
  \item Any direct-matter coupling for $\Xi$, any observable fifth-force
        prediction, or any laboratory-scale consequence.
  \item Any mathematical uniqueness theorem for $V = \Xi\log\Xi$
        derived within the present setting. The Bialynicki-Birula--Mycielski
        uniqueness theorem \cite{BialynickiBirulaMycielski1976} is cited
        as structural motivation, not as a proof of uniqueness under
        the cosmological hypotheses.
  \item Any analysis of the contribution of $\Xi$ during the
        inflationary epoch; the slow-roll parameters during inflation
        and any consequent effects on the primordial power spectrum are
        outside the scope of the present late-time dark-energy analysis.
\end{itemize}

\section{Summary}

A real, positive, dimensionless scalar field $\Xi$ with self-interaction
$V = \Lambda^{4}\,\Xi\log\Xi$, minimally coupled to gravity through
\eqref{eq:action}, defines a minimal dark-energy model with:
\begin{enumerate}
  \item Exact analytic vacuum $\Xi_0 = e^{-1}$
        (coupling-independent, universal).
  \item Stable massive fluctuations $m_\Xi^{2} = \Lambda^{4} e / M_\Pl^{2}$,
        with $\Lambda$ numerically near the observed dark-energy scale
        when $m_\Xi \sim H_0$.
  \item Freezing quintessence with asymptotic limit $w \to -1$.
  \item Structural relation $V = -\Lambda^{4}\,H_{\Gibbs}$ between the
        potential and the per-bin Gibbs integrand: the vacuum
        $\Xi_0 = e^{-1}$ coincides with the maximizer of $H_{\Gibbs}$
        on $(0,1]$. The connection is structural, not operational; see
        \S\ref{sec:entropy}.
  \item No observable fifth force is expected in the minimal theory.
  \item An illustrative consistency check against the DESI DR1 CPL
        central values: with $\Lambda \approx 1.7$ meV and the
        documented initial conditions, the model produces
        $(w_0, w_a)$ values within $\sim 1.2\sigma$ of the DR1 central
        values under the uncorrelated Gaussian-on-summary approximation,
        and a few-$\sigma$ from those central values under the
        correlated metric ($\rho \approx -0.9$). Both metrics quantify
        proximity to the published $(w_0, w_a)$ marginal Gaussian, not
        goodness-of-fit to the underlying data; the full DR2 +
        SN + CMB joint-likelihood analysis is deferred to a companion
        paper.
\end{enumerate}
The theory is falsifiable: a robust preference for phantom behavior on
the rolling branch, a non-monotone $w(z)$ trajectory, or failure of the
asymptotic approach to $-1$ would each rule out the minimal model at the
precision of Stage-IV surveys.

\section*{Acknowledgments}

The authors thank colleagues for discussions of the model's structural
properties. Verification scripts (\texttt{proof\_xi\_canonical.py}, 22
tests, all passing) and the DESI integration script
(\texttt{desi\_xi\_optimize\_v2.py}) are available as electronic
supplementary material. All errors are the authors'.

\begin{thebibliography}{99}

\bibitem{BialynickiBirulaMycielski1976}
I.~Bialynicki-Birula and J.~Mycielski,
\emph{Nonlinear wave mechanics},
Ann. Phys. (N.Y.) \textbf{100}, 62--93 (1976).
DOI: 10.1016/0003-4916(76)90057-9.

\bibitem{RatraPeebles1988}
B.~Ratra and P.~J.~E.~Peebles,
\emph{Cosmological consequences of a rolling homogeneous scalar field},
Phys. Rev. D \textbf{37}, 3406 (1988).

\bibitem{Wetterich1988}
C.~Wetterich,
\emph{Cosmology and the fate of dilatation symmetry},
Nucl. Phys. B \textbf{302}, 668 (1988).

\bibitem{Frieman1995}
J.~A.~Frieman, C.~T.~Hill, A.~Stebbins, and I.~Waga,
\emph{Cosmology with ultralight pseudo-Nambu-Goldstone bosons},
Phys. Rev. Lett. \textbf{75}, 2077 (1995).

\bibitem{CPL2001}
M.~Chevallier and D.~Polarski,
\emph{Accelerating universes with scaling dark matter},
Int. J. Mod. Phys. D \textbf{10}, 213 (2001).

\bibitem{BarrowParsons1995}
J.~D.~Barrow and P.~Parsons,
\emph{Inflationary models with logarithmic potentials},
Phys. Rev. D \textbf{52}, 5576 (1995); arXiv:astro-ph/9506049.

\bibitem{ColemanWeinberg1973}
S.~Coleman and E.~Weinberg,
\emph{Radiative corrections as the origin of spontaneous symmetry breaking},
Phys. Rev. D \textbf{7}, 1888 (1973).

\bibitem{Rosen1969}
G.~Rosen,
\emph{Dilatation covariance and exact solutions in local relativistic field theories},
Phys. Rev. \textbf{183}, 1186 (1969).

\bibitem{CazenaveHaraux1980}
T.~Cazenave and A.~Haraux,
\emph{Equations d'\'evolution avec non lin\'earit\'e logarithmique},
Ann. Fac. Sci. Toulouse \textbf{2}, 21--51 (1980).

\bibitem{HoeghKrohn1971}
R.~H\o egh-Krohn,
\emph{A general class of quantum fields without cut-offs in two space-time dimensions},
Commun. Math. Phys. \textbf{21}, 244--255 (1971).

\bibitem{Planck2020}
Planck Collaboration,
\emph{Planck 2018 results. VI. Cosmological parameters},
Astron. Astrophys. \textbf{641}, A6 (2020).

\bibitem{DESI2024BAO}
DESI Collaboration,
\emph{DESI 2024 III: Baryon acoustic oscillations from galaxies and quasars},
arXiv:2404.03000 (2024).

\bibitem{DESI2024VI}
DESI Collaboration,
\emph{DESI 2024 VI: Cosmological constraints from the measurements of baryon acoustic oscillations},
arXiv:2404.03002 (2024).

\bibitem{DESI2025DR2}
DESI Collaboration,
\emph{DESI DR2 results. II. Measurements of baryon acoustic oscillations and cosmological constraints},
arXiv:2503.14738; Phys. Rev. D \textbf{112}, 083515 (2025).

\bibitem{Wang2026QReconstruction}
S.~Wang, T.-N.~Li, T.~Liu, and G.-H.~Du,
\emph{Model-independent reconstruction of quintessence potential and kinetic energy from DESI DR2 and Pantheon+ supernovae},
arXiv:2603.21125 (2026).

\bibitem{Adil2026Analytic}
S.~A.~Adil, M.~A.~Zapata, \"{O}.~Akarsu, and J.~A.~Vazquez,
\emph{Background-level reconstruction of scalar-field potentials from dark-energy histories and comparison with analytic potential families},
arXiv:2603.14693 (2026).

\bibitem{AdamEtAl2025}
H.~Adam, M.~P.~Hertzberg, D.~Jim\'enez-Aguilar, and I.~Khan,
\emph{Comparing Minimal and Non-Minimal Quintessence Models to 2025 DESI Data},
arXiv:2509.13302 (2025).

\bibitem{Turyshev2026Review}
S.~G.~Turyshev,
\emph{Dark Energy After DESI DR2: Observational Status, Reconstructions, and Physical Models},
arXiv:2602.05368 (2026).

\bibitem{ShajibFrieman2025}
A.~J.~Shajib and J.~A.~Frieman,
\emph{Scalar field dark energy models: current and forecast constraints},
arXiv:2502.06929; Phys. Rev. D \textbf{112}, 063508 (2025).

\bibitem{Thompson2019}
R.~I.~Thompson,
\emph{Beta function quintessence cosmological parameters and fundamental
constants -- II. Exponential and logarithmic dark energy potentials},
Mon. Not. Roy. Astron. Soc. \textbf{482}, 5448 (2019); arXiv:1811.03164.

\bibitem{KouwnOh2012}
S.~Kouwn and P.~Oh,
\emph{Dark energy with logarithmic cosmological fluid},
J. Korean Phys. Soc. \textbf{65}, 814 (2014); arXiv:1201.4544.

\bibitem{Chavanis2015}
P.-H.~Chavanis,
\emph{Logotropic dark fluid as a unification of dark matter and dark energy},
Phys. Rev. D \textbf{92}, 103004 (2015); arXiv:1505.00034.

\bibitem{Oikonomou2019}
V.~K.~Oikonomou,
\emph{Generalized logarithmic equation of state in classical and loop
quantum cosmology dark energy--dark matter coupled systems},
Annals of Physics \textbf{409}, 167912 (2019); arXiv:1907.02600.

\bibitem{GW170817}
B.~P.~Abbott \emph{et al.} (LIGO Scientific Collaboration and Virgo Collaboration),
\emph{Gravitational waves and gamma-rays from a binary neutron star merger:
GW170817 and GRB 170817A},
Astrophys. J. Lett. \textbf{848}, L13 (2017).

\bibitem{Verlinde2011}
E.~Verlinde,
\emph{On the origin of gravity and the laws of Newton},
JHEP \textbf{04}, 029 (2011).

\bibitem{CHLZ2012}
A.~Sheykhi \emph{et al.},
\emph{Entropy-corrected holographic dark energy and interacting quintessence reconstruction},
Astrophys. Space Sci. \textbf{342}, 93 (2012).

\bibitem{Ensslin2013}
T.~A.~En{\ss}lin,
\emph{Information theory for fields},
Annalen der Physik \textbf{525}, 895 (2013); arXiv:1301.2556.

\bibitem{Caticha2012}
A.~Caticha,
\emph{Entropic inference and the foundations of physics},
EBEB 2012 monograph (2012); arXiv:1412.5629.

\bibitem{EotWash}
D.~J.~Kapner \emph{et al.},
\emph{Tests of the gravitational inverse-square law below the dark-energy length scale},
Phys. Rev. Lett. \textbf{98}, 021101 (2007).

\bibitem{MICROSCOPE}
P.~Touboul \emph{et al.},
\emph{MICROSCOPE mission: final results of the test of the equivalence principle},
Phys. Rev. Lett. \textbf{129}, 121102 (2022).

\bibitem{SandersGish2026Sigma}
B.~R.~Sanders and M.~Gish,
\emph{Non-Associativity Decay in Binary Composition Tables over $\mathbb{Z}/N\mathbb{Z}$},
Preprint, 2026. Companion paper; archived in shared deposit DOI 10.5281/zenodo.18852047.

\bibitem{SandersGish2026FourCore}
B.~R.~Sanders and M.~Gish,
\emph{Joint Closure, Per-Coordinate Fuse Data, and a Closed-Form Algebraic
Attractor of Two Commutative Binary Operations on $\mathbb{Z}/10\mathbb{Z}$},
Preprint, 2026. Companion paper; archived in shared deposit DOI 10.5281/zenodo.18852047.

\end{thebibliography}

\end{document}